% Part: counterfactuals
% Chapter: minimal-change-semantics
% Section: transitivity

\documentclass[../../../include/open-logic-section]{subfiles}

\begin{document}

\olfileid[zh]{cnt}{min}{cpo}

\olsection{逆否命题}

实质条件句和严格条件句都等价于其逆否命题。反事实条件句则不然。下面这个例子来自 Kratzer：
\begin{quote}
  如果歌德不是在 1832 年去世的，那么他（现在依然）已经死了。

  如果歌德现在还没有死，那么他就是在 1832 年去世的。
\end{quote}
第一句为真：人活不了几百年。第二句显然为假：如果歌德现在还没有死，那么他仍然在世，因此他不可能是在 1832 年去世的。

\begin{ex}\ollabel{ex:contraposition-counterex}
  球面语义使得逆否命题无效，即我们有 $p
  \cif q \Entails/ \lnot q \cif \lnot p$。把 $p$ 想作「歌德
  不是在 1832 年去世的」，把 $q$ 想作「歌德现在就死了」。我们可以在模型 $\mModel{M_1} = \tuple{W, O, V}$ 中刻画这一点，其中 $W =
  \{w, w_1, w_2\}$，$O = \{\{w\}, \{w, w_1\}, \{w, w_1, w_2\}\}$，
  $V(p) = \{w_1, w_2\}$ 且 $V(q) = \{w, w_1\}$。因此 $w$ 是歌德在 1832 年去世并且至今仍死的现实世界；$w_1$ 是歌德在（比如说）1833 年去世并且至今仍死的（近）世界；而 $w_2$ 是歌德至今仍在世的（远）世界。
\begin{figure}
\begin{center}
\begin{tikzpicture}[scale=.6,layerwidth=1.5]\tiny
  \spheresystem{3}
  \begin{scope}
    \clip \propositionplot{0}{.8}{50}{4.5} -- (4.5,-4.5) -- (-4.5,-4.5) -- (-4.5,4.5) -- (4.5,4.5);
    \spherelayer{3}
  \end{scope}
  \propositionintersect{0}{2}{110}{5.5}
  \proposition{0}{.8}{50}{4.5}
  \path (0,0) node[world] {$w$};
  \spherepos{0}{2}{node[world] {$w_1$}}
  \spherepos{-50}{3}{node[world] {$w_2$}}
  \spherepos{28}{3}{node {$q$}}
  \spherepos{42}{3}{node {$\lnot q$}}
  \spherepos{-55}{4}{node {$p$}}
  \spherepos{-67}{4}{node {$\lnot p$}}
\end{tikzpicture}
\caption{逆否命题的反例}
\ollabel{fig:contraposition}
\end{center}
\end{figure}
存在一个 $p$-容纳球面 $S =
\{w, w_1\}$，并且 $p \lif q$ 在其中所有世界处都为真，因此 $\mSat{M}{p
  \cif q}[w]$。然而，$\lnot q$-容纳球面 $\{w, w_1,
w_2\}$ 含有一个世界，即~$w_2$，在其中 $q$ 为假而 $p$ 为
真，因此 $\mSat/{M}{\lnot q \lif \lnot p}[w_2]$。
\end{ex}

\end{document}
